% generated by scripts/make_cmame_material.py from the records; do not edit
\begingroup\small
\begin{longtable}{@{}>{\raggedright\arraybackslash}p{0.145\linewidth}>{\raggedright\arraybackslash}p{0.405\linewidth}>{\raggedright\arraybackslash}p{0.21\linewidth}>{\raggedright\arraybackslash}p{0.12\linewidth}@{}}
\caption{Every pre-registered criterion of the runs reported here, with its outcome.}\label{tab:criteria}\\
\toprule
Run & Criterion & Measured & Outcome \\
\midrule
\endfirsthead
\caption[]{Every pre-registered criterion of the runs reported here, with its outcome. (continued)}\\
\toprule
Run & Criterion & Measured & Outcome \\
\midrule
\endhead
\bottomrule
\endfoot
July run (2D), 16 July 2026 & G1$'$ (a), energy-term sub-experiment (separately trained networks): supervised + energy term at least 3 times better than the zero field in displacement, and better than the two naive baselines and the plain polynomial, fitted on 1,024 labelled instances, at every budget & 4.5--6.1 times; baselines beaten at every budget & met \\
 & G1$'$ (b), \cref{tab:2djuly,tab:labeleffjuly}, 64 labels: the energy term lowers the relative energy gap by $\geq$ 50\% and the von Mises error by $\geq$ 25\% & 78.7\%, 65.4\% & met \\
 & G1$'$ (c), \cref{tab:2djuly,tab:labeleffjuly}, 64 labels: label-free, then supervised beats supervised by $\geq$ 10\% in relative energy gap or $\geq$ 5\% in displacement & 92.6\%, 2.0\% & met \\
 & Gate G1$'$ = (a) and (b) and (c) &  & GO \\
 & K1, energy-term sub-experiment: the better of the fixed and the gradient-scaled energy term lowers the relative energy gap by $<$ 25\% at every budget & 98.0\% / 94.4\% / 90.6\% / 74.0\% & not triggered \\
 & K2, \cref{tab:2djuly}: label-free displacement error $>$ 30\% above supervised, 1,024 labels & +3.6\% & not triggered \\
 & C1-advantage, \cref{tab:labeleffjuly}: label-free relative energy-gap advantage over supervised $<$ 40\% at every budget & 98.7\% / 97.0\% / 93.5\% / 80.4\% & not triggered \\
 & K4: a latent-space regulariser raises the standardised effective rank of the latent vectors by $\leq$ 1.5 times & 1.70 times & not triggered \\
 & K5: conjugate gradients from the prediction of a label-free network (64 validation instances) save $<$ 20\% of the iterations from zero & $-$0.9\% & triggered \\
 & K6: a cross-resolution invariance term shrinks the transfer gap by $<$ 10\% at coarsening 2.5 & $-$23.7\% & triggered \\
 & E6 (alignment): within-geometry rank correlation of latent and solution distances $<$ 0.3 & 0.730 & not triggered \\
 & C5: any pre-registered check of the conditioning bound, the modewise identity or the conjugate-gradient bound violated & none & not triggered \\
\midrule
Joint-embedding run (2D), 31 July 2026 & K3: a self-supervised joint-embedding term added to label-free pretraining changes the displacement error and the relative energy gap of the transformer, fine-tuned with 16, 64 and 256 labels, by less than 3\% at every budget (if triggered, the term is dropped) & change with the term, 16 / 64 / 256 labels: displacement $-$13.5\% / $-$0.5\% / +2.6\%; relative energy gap $-$48.4\% / +92.5\% / $-$16.4\% (positive: more accurate) & not triggered \\
\midrule
Initial three-dimensional run, 21 September 2026 & Gate G2 and kill conditions as registered; the label-free objective was defective (deviation D14) & (a), (b), (c) not met; KP1, KP2, KP4 triggered & NO-GO \\
Pilot of the amendment (3D), 22 September 2026 & Label-free displacement error after 20 epochs $<$ 0.90 & 0.309 & met \\
Three-dimensional comparison, 28 September 2026 & G2 (a), budgets $\geq$ 64 (amendment), \cref{tab:labeleff}: supervised + energy term at least 3 times better than the zero field in displacement, and better than both naive baselines & 3.9 / 12.0 / 21.3 times; naive beaten & met \\
 & G2 (a), every budget (the registered form, reported for reference) & 2.71 times at 16 labels & not met \\
 & G2 (b), \cref{tab:labeleff,tab:3d}: label-free displacement error within +10\% of supervised at 1,024 labels, and relative energy-gap advantage $\geq$ 40\% at every budget & $-$36.3\%; 99.95\% / 99.85\% / 98.86\% / 99.97\% & met \\
 & G2 (c), \cref{tab:transfer}: finer mesh, zero-shot displacement error $\leq$ 1.25 times in-band, and better than both naive baselines & 8.62 times; best naive 0.401 against 0.258 & not met \\
 & Gate G2 = (a) and ((b) or (c)), amended &  & GO \\
 & Gate G2, registered form (reference) &  & NO-GO \\
 & KP1: supervised (1,024 labels) better than label-free by $>$ 10\% in displacement & $-$36.3\% & not triggered \\
 & KP2: relative energy-gap advantage $<$ 40\% at any budget & minimum 98.9\% & not triggered \\
 & KP3: the energy term lowers the relative energy gap by $<$ 25\% at every budget & +96.6\% / +85.7\% / $-$20.7\% / +99.6\% & not triggered \\
 & KP4: finer mesh, zero-shot displacement error $>$ 1.5 times in-band, or a naive baseline better & 8.62 times & triggered \\
 & KP5: any pre-registered check of the conditioning bound, the modewise identity or the conjugate-gradient bound violated & none & not triggered \\
 & KP6 (alignment): within-geometry rank correlation of latent and solution distances $<$ 0.3 & 0.827 & not triggered \\
\midrule
Latent-\hspace{0pt}regularisation run (2D), 29 September 2026 & Latent regularisation, K1: relative change of the displacement error or of the relative energy gap beyond the noise guard & within the guard & not triggered \\
 & Latent regularisation, K2: the separation does not rise in any seed & rose in every seed & not triggered \\
 & Latent regularisation, GO: separation raised by $\geq$ 0.02 in every seed & +0.009 / +0.014 / +0.002 & NO-GO \\
\midrule
Bottleneck run (3D), 30 September 2026 & Token bottleneck, M = 512, K1: in-band relative energy gap or finer-mesh displacement error worse than the label-free transformer's beyond the noise guard & +280.9\% (guard 287.2\%), +79.1\% (guard 41.9\%) & triggered \\
 & Token bottleneck, M = 512, K2: finer-mesh training step longer than 2 s & 0.049 s & not triggered \\
 & Token bottleneck, M = 512: outcome &  & dropped \\
 & Token bottleneck, M = 1,024, K1: in-band relative energy gap or finer-mesh displacement error worse than the label-free transformer's beyond the noise guard & +183.3\% (guard 70.5\%), +122.0\% (guard 33.3\%) & triggered \\
 & Token bottleneck, M = 1,024, K2: finer-mesh training step longer than 2 s & 0.047 s & not triggered \\
 & Token bottleneck, M = 1,024: outcome &  & dropped \\
\midrule
Two-dimensional comparison, 9 October 2026 & H1, \cref{tab:2d,tab:cm2d}: label-free transformer lower than the supervised transformer ($\mathcal{L}_D$) in relative energy gap beyond the noise guard, 1,024 labels & $-$71.3\% (guard 12.4\%) & supported \\
 & H2a: stiffness-norm transformer ($\mathcal{L}_K$) lower than the supervised transformer ($\mathcal{L}_D$) in relative energy gap beyond the guard & $-$78.2\% (guard 16.4\%) & supported \\
 & H2b: stiffness-norm transformer lower than the supervised transformer in von Mises error beyond the guard & $-$54.9\% (guard 12.3\%) & supported \\
 & H3 (a reading, no criterion): stiffness-norm transformer against the label-free transformer in relative energy gap & $-$24.2\% (guard 38.1\%) & no difference shown \\
 & Reuse: the label-free networks of the latent-regularisation run reproduce that run's per-instance relative energy gaps and displacement errors (largest relative deviation at most $10^{-4}$) & largest relative deviation 0 & passed \\
\end{longtable}
\endgroup
\vspace{-\smallskipamount}
\noindent\begin{minipage}{\linewidth}\footnotesize\raggedright
Runs are named as in \cref{tab:runs} and dated by their records; the pilot of the amendment is the short run that the amendment required before the three-dimensional comparison (\ref{app:prereg}). Labels are those of the pre-registration documents, each of which numbers its own criteria. Gates (G1$'$, G2) decide whether the study proceeds to its next stage. Kill conditions withdraw a claim, or abandon a component, when triggered: in two dimensions K1--K6, C1-advantage (the claimed advantage in relative energy gap), C5 (the pre-registered checks of \ref{app:repro:checks}) and E6 (latent alignment); in three dimensions KP1--KP6, with KP5 the counterpart of C5. The latent-regularisation and bottleneck runs number their criteria afresh, with the meanings given here. A criterion reads the trained networks of the tables it names; the energy-term sub-experiment of the July run trained its own supervised networks with and without the energy term, with the same configuration. K4, E6 and KP6 concern latent representations, which this paper does not use. Standardised effective rank: the participation ratio of the eigenvalues of the covariance of the latent vectors, each dimension standardised. The 3D amendment, the deviations and the joint-embedding term are described below, the latent regularisation and the token bottleneck in \ref{app:e1}, the July run in \ref{app:july}.\par
Noise guard: a comparison of a network A with a reference B over three seeds each, with rel = mean(A)/mean(B) $-$ 1 and SE$_\mathrm{rel}$ = $(s_A^2/3 + s_B^2/3)^{1/2}$ divided by mean(B), where $s_A$ and $s_B$ are the sample standard deviations over the seeds, finds A lower beyond the guard when rel is below minus the guard, max(10\%, 2 SE$_\mathrm{rel}$), and higher beyond it when rel exceeds the guard; within it no difference is shown, which is not equivalence. The hypotheses H1--H3 of the 2D comparison use it (\cref{tab:cm2d}), and so do the latent-regularisation and bottleneck runs, with the label-free network as B; their pre-registrations call it the resolution. $\mathcal{L}_D$ and $\mathcal{L}_K$ are the losses \eqref{eq:ld} and \eqref{eq:lk}.
\end{minipage}
\par\medskip
